\documentclass[12pt,a4paper]{article}

\usepackage{xeCJK}
\usepackage{fontspec}
\usepackage{geometry}
\usepackage{titlesec}
\usepackage{enumitem}
\usepackage{xcolor}
\usepackage{hyperref}
\usepackage{amsmath}

\geometry{margin=2.5cm}

% 字体设置
\IfFontExistsTF{Noto Serif CJK SC}{
  \setCJKmainfont{Noto Serif CJK SC}
  \setCJKsansfont{Noto Sans CJK SC}
}{
  \setCJKmainfont{STSong}
  \setCJKsansfont{STHeiti}
}

% 标题格式
\titleformat{\section}{\large\bfseries\color{darkred}}{\thesection}{1em}{}
\titleformat{\subsection}{\normalsize\bfseries}{}{0em}{}

\definecolor{darkred}{RGB}{139,0,0}
\definecolor{darkgreen}{RGB}{0,100,0}

\setlength{\parskip}{0.5em}
\setlength{\parindent}{0em}

\title{\textbf{答辩讲稿}\\[0.5em]\large 同类产品评论展示策略研究：基于消费者学习的分析}
\author{答辩人：焦洋 \quad 指导教师：张瑞洁}
\date{预计时长：20分钟}

\begin{document}

\maketitle
\thispagestyle{empty}

\vspace{1em}
\hrule
\vspace{2em}

\section*{第1页：标题页（30秒）}

各位老师好，我是焦洋，我的毕业论文题目是《同类产品评论展示策略研究：基于消费者学习的分析》，指导教师是张瑞洁老师。下面我将用大约20分钟的时间向各位老师汇报我的研究工作。

\section*{第2页：目录（20秒）}

我的汇报将按照以下顺序展开：首先介绍研究背景与动机，然后是文献综述，接着是模型设定，之后分别分析分开展示和聚合展示两种策略，然后进行策略比较和数值分析，最后是福利分析、模型扩展以及结论。

\section*{第3页：研究背景（1分钟）}

首先是研究背景。

在电子商务环境中，消费者面临严重的信息不对称问题——他们在购买前无法直接体验产品质量。在线评论因此成为消费者获取产品信息的关键来源。

那么，平台应该如何设计评论展示系统呢？这是我们研究的核心问题。

我们关注的是一个典型场景：\textbf{垂直差异化产品}。比如同品牌不同型号的手机，iPhone 15和iPhone 16；或者同类别不同配置的产品，比如基础款和高端款。这些产品之间存在质量差异，但消费者在购买前很难准确判断。

\section*{第4页：两种评论展示策略（1分钟）}

针对这类产品，平台主要有两种评论展示策略。

第一种是\textbf{分开展示策略}，也就是各产品的评论独立展示。消费者查看某个产品时，只能看到该产品自己的评论。这是大多数电商平台的默认方式。

第二种是\textbf{聚合展示策略}，也就是将所有同类产品的评论合并展示。消费者看到的是一个综合评分。部分平台在展示"系列产品"评分时会采用这种方式。

这两种策略各有特点：分开展示信息更精准但评论数量可能较少；聚合展示评论数量更多但信息可能被稀释。那么平台应该选择哪种策略呢？

\section*{第5页：研究问题（40秒）}

基于以上背景，我们的核心研究问题是：对于销售垂直差异化产品的平台，应该选择哪种评论展示策略？

具体来说，我们关注四个子问题：
\begin{enumerate}[itemsep=0pt]
    \item 两种策略下消费者的学习行为有何不同？
    \item 策略选择如何影响平台利润？
    \item 策略选择对消费者福利有何影响？
    \item 评论精度的非对称性如何影响策略比较？
\end{enumerate}

\section*{第6页：文献综述（1分钟）}

本研究与四个领域的文献相关。

首先是\textbf{信息设计与贝叶斯说服}领域，Kamenica和Gentzkow 2011年的开创性工作为我们提供了理论框架基础。

第二是\textbf{在线评论与消费者学习}领域，Acemoglu等人2022年在Econometrica上的工作研究了评论学习机制。

第三是\textbf{平台评论系统设计}领域，这是与本文最直接相关的研究，包括Xiao等人和Shi等人的工作。

第四是\textbf{多产品定价}领域，为我们的产品线背景提供了理论支撑。

本文的主要贡献在于：首次研究同类产品间评论聚合与分开展示的策略选择问题，并且将信息的内生生成过程纳入信息设计框架。

\section*{第7页：模型概述（1分钟）}

下面介绍模型设定。

我们考虑一个垄断平台销售两种产品：产品B是基础版，质量服从均值为 $q_L$ 的正态分布；产品G是高级版，质量服从均值为 $q_H$ 的正态分布，其中 $q_H > q_L$。两种产品共享相同的先验精度 $\tau$。

模型采用两期结构。第一期，消费者基于先验信念购买产品，购买者中有一定比例会发布评论。第二期，新一批消费者观察评论后更新信念，然后做出购买决策。

\section*{第8页：需求与评论机制（50秒）}

需求函数采用标准的线性形式：购买量等于市场规模乘以感知价值减去价格敏感度乘以价格。

评论生成方面，我们假设购买者中有比例 $\theta$ 会发布评论，评论内容反映消费者的感知质量，服从以真实质量为均值的正态分布。

评分信号方面，在分开展示下，每种产品有独立的评分信号；在聚合展示下，平台将两种产品的评论合并为一个加权平均的综合评分。

\section*{第9页：基本假设（40秒）}

模型有三个基本假设。

假设一是策略承诺：平台在定价之前承诺评论展示策略，这反映了展示格式作为平台设计的制度性约束。

假设二是评论真实性：所有评论均为真实体验反馈，不存在虚假评论。

假设三是理性贝叶斯学习：消费者能够正确地根据所观察到的信息更新质量信念。

\section*{第10页：平台利润函数（1分钟）}

下面推导平台的利润函数，这是理解策略比较的关键。

首先看第一期利润。由需求函数可以反解出价格，代入后得到第一期利润是购买量 $k$ 的二次函数。第一期消费者基于先验信念购买，感知价值就是质量的先验均值。

然后看第二期。第二期消费者观察评论后更新信念，感知价值变成后验均值 $v_\delta$。平台选择价格最大化利润，求解一阶条件得到最优价格是感知价值除以 $2h$，对应的最优利润是感知价值平方除以 $4h$，再乘以市场规模。

注意这里感知价值可能为负，所以要取 max 和零。

\section*{第11页：利润与后验方差的关系（1分钟）}

这一页是理解核心结论的关键。

第二期的期望利润取决于 $\E[\max\{0, v_\delta\}^2]$。由于后验均值 $v_\delta$ 服从正态分布，均值是真实质量均值（大于零），方差是 $\sigma^2_\delta$。

这里有一个重要的引理：对于均值为正的正态随机变量，它的期望平方正部关于方差是\textbf{递增的}。

经济含义是什么呢？后验方差越大，说明评分对消费者信念的影响越大，信息越有价值。平台可以根据已揭示的质量信息调整价格——质量信号好时提高价格，质量信号差时降低价格。这种状态依存定价带来的期望收益是正的。

由于分开展示的后验方差大于聚合展示，这直接推出分开展示的利润更高。这就是我们核心结论的数学基础。

\section*{第12页：分开展示的贝叶斯更新（50秒）}

现在分析分开展示策略。

在分开展示下，消费者观察独立的评分信号，根据贝叶斯法则进行更新。后验分布仍然是正态分布，后验均值是先验均值与评分的精度加权平均。

从平台的事前视角看，后验均值本身是一个随机变量。它的期望等于真实质量均值，方差为 $\sigma_{S\delta}^2$，具体公式如幻灯片所示。

分开展示的一个关键特征是：两种产品的定价决策\textbf{相互独立}。因为每种产品的评分只取决于该产品自身的购买量。

\section*{第11页：聚合展示的贝叶斯更新（1分钟）}

聚合展示下的贝叶斯更新更加复杂。

消费者只观察到一个综合评分 $r_I$，要推断某一产品的质量，需要对另一产品的未知质量进行积分。后验均值的表达式如幻灯片所示，结构比分开展示复杂得多。

聚合展示的关键特征是：两种产品的定价决策\textbf{相互耦合}。增加一种产品的销量会影响消费者对两种产品的学习。

\section*{第12页：聚合展示降低信息精度（1分钟）}

这里是我们的第一个重要结论，命题5.2：\textbf{信息稀释效应}。

对于任意给定的购买量组合，聚合展示下消费者信念的方差不超过分开展示。也就是说，$\sigma_{I\delta}^2 \leq \sigma_{S\delta}^2$。

这个结论的经济含义是：聚合展示混合了不同质量产品的评论，引入了\textbf{评论来源不确定性}。消费者无法区分综合评分中的信息来自哪种产品。

这里需要注意的是，后验方差更低意味着评分对信念的影响更小，也就是说消费者从评论中学到的更少。所以信息稀释效应实际上降低了学习效率。

\section*{第13页：核心结果一（1分钟）}

现在进入策略比较部分，这是本文的核心。

命题5.3表明：对于任意给定的购买量，只要两种产品都有销售，分开展示的期望总利润\textbf{严格高于}聚合展示。

证明思路是这样的：第一期利润在给定相同购买量下是相同的，利润差异完全来自第二期。对于均值为正的正态随机变量，期望平方的正部是关于方差严格递增的。分开展示方差更大，因此第二期利润更高。

\section*{第14页：核心结果二（1分钟）}

命题5.4给出了一个精确的量化结果：在对称情形下，也就是两种产品购买量相同时，分开展示对第二期利润的期望贡献恰好是聚合展示的\textbf{两倍}。信息效率比恰好是2:1。

这个结果的经济解释是：虽然聚合展示将样本量从 $k$ 增加到了 $2k$，但由于来源不确定性，消费者无法有效利用这些额外的样本。净效果是有效信息量减半。

值得强调的是，这个2:1的比例与质量差距 $\Delta q$ 无关。即使两种产品质量非常接近，聚合展示仍然会损失一半的信息效率。

\section*{第15页：核心结果三（1分钟）}

命题5.5是基于数值模拟的经验性结论：在我们模型设定下，对于所有测试的参数配置，分开展示在最优决策下均产生严格更高的期望利润。

为什么聚合展示无法逆转劣势呢？有三个原因：
\begin{enumerate}[itemsep=0pt]
    \item 信息价值为正：精确的质量信号使平台能够进行状态依存定价，带来正的期望收益
    \item 样本规模优势不足：来源不确定性的成本超过了样本合并的收益
    \item 交叉效应有限：最优购买量的差异只有5-7\%，不足以改变策略排序
\end{enumerate}

\section*{第16页：数值模拟设定（40秒）}

数值模拟采用的基准参数如表所示。低质量产品均值为1，高质量产品均值为2，先验精度、感知精度、价格敏感度和市场规模都设为1，评论生成概率为0.5。

在基准参数下，分开展示的最优利润为2.6355，聚合展示为2.5741，利润差为0.0613。

\section*{第17页：敏感性分析（50秒）}

敏感性分析检验了结论的稳健性。图中显示了不同参数对利润差的影响。

最关键的发现是：利润差始终为正，也就是说在所有测试的参数范围内，分开展示始终占优。

其中，先验精度 $\tau$ 的影响最显著。当 $\tau$ 较低，也就是消费者先验不确定性较高时，评论信息的价值更大，分开展示的优势也更明显。

质量差距 $\Delta q$ 的影响相对有限，评论生成率 $\theta$ 越高，分开展示的优势越明显。

\section*{第18页：福利分析（50秒）}

福利分析得到了一个重要结论，命题5.6：分开展示策略同时最大化平台利润、消费者剩余和社会总福利。

这意味着平台利润最大化与社会福利最大化指向同一策略，\textbf{不存在利润-福利冲突}。

这与信息设计文献中常见的权衡不同。原因在于我们模型的特殊结构：在线性需求设定下，消费者剩余与平台利润保持固定比例，因此两个目标是一致的。

\section*{第19页：扩展——评论精度非对称性（1分钟）}

最后介绍一个重要的模型扩展：评论精度非对称性。

在现实中，不同产品的评论质量可能不同。高质量产品的用户可能写出更详细、更精确的评论；不同产品的用户群体评论能力也可能不同。

我们放松了基准模型中两种产品评论精度相同的假设，定义精度比 $\rho = \zeta_G / \zeta_B$。

命题6.1表明：当精度对称时，信息效率比恰好是2；当精度不对称时，信息效率比严格大于2。也就是说，精度非对称性会\textbf{放大}分开展示的优势。

\section*{第20页：扩展数值验证（40秒）}

数值结果验证了这一理论预测。表格显示了不同精度比下的利润比较。

当 $\rho = 1$ 时，信息比恰好是2，与我们的理论一致。当 $\rho$ 偏离1时，无论是大于1还是小于1，信息比都严格大于2，分开展示的优势都会扩大。

命题6.2总结道：基准模型中"分开展示普遍占优"的结论对评论精度的非对称性具有稳健性。精度非对称性是分开展示优势的放大因素，而非逆转因素。

\section*{第21页：主要结论（1分钟）}

总结本文的主要结论。

第一，聚合展示降低信息精度。它引入评论来源不确定性，降低消费者学习效率。

第二，分开展示普遍占优。对称情形下信息效率是2:1，在所有测试参数下分开展示均占优。

第三，利润与福利目标一致。分开展示同时最大化利润、消费者剩余和社会福利，不存在利润-福利冲突。

第四，评论精度非对称性强化结论。精度差异越大，分开展示的优势越明显。

\section*{第22页：管理启示与未来方向（50秒）}

基于以上结论，我们可以得到几点管理启示。

对于垂直差异化产品组合，平台应优先采用分开展示。高质量评论是信息资产，不应与低质量评论混合稀释。平台可以通过评论标签、评论者等级等机制进一步提升评论精度的差异化价值。

未来研究可以从以下方向扩展：模型方面可以考虑多产品、动态评论积累和竞争平台；机制方面可以引入消费者异质性学习和评论操纵成本；实证方面可以利用平台数据检验理论预测。

\section*{第23页：参考文献（10秒）}

这里列出了本研究引用的主要参考文献。

\section*{第24页：致谢（20秒）}

以上就是我的汇报内容。感谢张瑞洁老师在论文写作过程中的悉心指导，也感谢各位老师的聆听。

欢迎各位老师提问和指正，谢谢！

\newpage
\section*{备用问答要点}

\subsection*{Q: 为什么用正态分布？}
A: 正态分布的共轭性质使得贝叶斯更新有闭式解，便于理论分析。同时正态分布的对称性简化了模型，使我们能够聚焦于评论展示策略的核心权衡。

\subsection*{Q: 2:1的信息效率比是怎么来的？}
A: 在对称情形下，分开展示的后验方差是 $k\theta\zeta / (\tau^2 + k\tau\theta\zeta)$，聚合展示是这个值的一半。虽然聚合展示样本量翻倍，但来源不确定性恰好抵消了这个收益。

\subsection*{Q: 什么情况下聚合展示可能占优？}
A: 在我们的基准模型下不会。但在扩展方向中，如果考虑消费者偏好评论数量（有限理性）、平台面临信誉成本、或者存在产品多样性的社会价值，聚合展示可能在某些条件下占优。

\subsection*{Q: 线性需求假设的局限性？}
A: 线性需求使得消费者剩余与平台利润保持固定比例，这导致利润-福利目标一致。在更一般的需求函数下，可能出现两个目标的冲突。

\subsection*{Q: 如何验证理论预测？}
A: 可以利用平台A/B测试数据，比较同一产品系列在不同展示策略下的销售和价格反应。也可以寻找平台评论合并规则变化前后的自然实验。

\end{document}
